\documentclass[journal]{IEEEtran}
\usepackage{cite}
\usepackage{amsmath,amssymb,amsfonts}
\usepackage{algorithmic}
\usepackage{graphicx}
\usepackage{textcomp}
\usepackage{xcolor}
\usepackage{booktabs}
\usepackage{url}
\usepackage{microtype}

\begin{document}

\title{TRUEGUARD: A Lightweight Real-Time Framework for Hallucination Detection and Mitigation in Large Language Models}

\author{Yogesh~Ravi~M and~T.~S.~Pradeep~Kumar%
\thanks{Yogesh Ravi M is with the School of Computer Science and Engineering, Vellore Institute of Technology, Chennai 600127, India (e-mail: yogeshravi.m2022@vitstudent.ac.in).}%
\thanks{T. S. Pradeep Kumar is with the Faculty of Computer Science and Engineering, Vellore Institute of Technology, Chennai 600127, India (e-mail: pradeepkumar.ts@vit.ac.in).}%
}

\markboth{IEEE Transactions on Artificial Intelligence,~Vol.~XX, No.~X,~2026}%
{Ravi and Pradeep Kumar: TRUEGUARD: Real-Time Hallucination Detection and Mitigation}

\maketitle

\begin{abstract}
Large Language Models (LLMs) frequently generate fluent yet factually unfounded statements---a phenomenon known as hallucination. Existing detection approaches generally mandate either multi-sample semantic consistency checks (which require 5--10 stochastic forward passes and introduce $\sim$142.3~ms/token overhead) or isolated single-signal white-box metrics that fail to capture heterogeneous hallucination modes. In this paper, we propose TRUEGUARD, an inline, multi-signal framework that monitors autoregressive generation in a single forward pass without disrupting standard decoding. TRUEGUARD extracts four complementary intrinsic signals at each decoding step: Internal State Probes (ISP), Attention Anomaly Score (AAS), Entropy-Energy Monitor (EEM), and Distribution Shift Tracker (DST). These signals are dynamically standardized via online $Z$-score normalization and fused through a learned logistic regression layer to produce a calibrated token-level risk score $R(t)$. When the sequence risk exceeds a calibrated safety threshold ($\tau = 0.50$), TRUEGUARD triggers an active closed-loop mitigation pipeline incorporating targeted Retrieval-Augmented Grounding (RAG) and calibrated abstention. Evaluated across five benchmark distributions (TruthfulQA, HaluEval-QA, HaluEval-Summ, HaluEval-Dial, and HELM), TRUEGUARD achieves an average AUROC of $0.8821 \pm 0.0170$ (with $0.8533 \pm 0.0533$ on TruthfulQA) while imposing merely $12.4 \pm 0.5$~ms/token latency overhead---an 11.5$\times$ speedup over Semantic Entropy. In human-annotated mitigation evaluations ($n = 300$, Cohen's $\kappa = 0.82$), TRUEGUARD achieves a 63.8\% relative reduction in hallucination occurrence while maintaining an 88.0\% query answer coverage. Complete reproducible code, Colab notebooks, and evaluation harnesses are available at https://github.com/yogeshravi19/trueguard.
\end{abstract}

\begin{IEEEkeywords}
Large Language Models, Hallucination Detection, Trustworthy AI, Real-Time Inference, Uncertainty Quantification, Active Mitigation, Explainable AI.
\end{IEEEkeywords}

\section{Introduction}
\IEEEPARstart{L}{arge} Language Models (LLMs) such as GPT-4, LLaMA, and Mistral have demonstrated remarkable empirical capabilities in language generation, reasoning, and synthesis \cite{huang2024trustllm, zhang2025siren}. Despite their widespread adoption across diverse software environments, their deployment in safety-critical domains---including healthcare, legal informatics, and automated financial analysis---remains severely constrained by their propensity to hallucinate \cite{ji2023survey}. Hallucinations occur when an LLM produces outputs that appear syntactically coherent and contextually plausible but are factually false, unfaithful to the prompt context, or logically contradictory \cite{min2023factscore}.

Prior research in hallucination detection has largely pursued two contrasting design philosophies:
\begin{enumerate}
    \item \textbf{Black-Box Multi-Sample Consistency Checks:} Methods such as SelfCheckGPT \cite{manakul2023selfcheckgpt} and Semantic Entropy \cite{kossen2024semantic} sample multiple stochastic completions per prompt and evaluate their semantic variance. While these methods achieve reliable detection accuracy, generating multiple responses consumes significant memory bandwidth and compute time (incurring $\sim$142.3~ms/token overhead), making them impractical for low-latency streaming applications.
    \item \textbf{Single-Signal White-Box Probes:} To bypass multi-pass latency, recent works monitor internal model activations in a single forward pass. Representative examples include internal representation probes \cite{su2024mind, zhang2025mhad}, attention pattern tracking \cite{hajji2025map}, semantic energy logits analysis \cite{ma2025semantic}, and distribution shift trackers like HaluShift \cite{dasgupta2025hallushift}. However, individual internal signals typically capture only specific facets of uncertainty. As empirical studies demonstrate, attention dispersion effectively flags intrinsic contradictions against context tokens but often misses extrinsic fabrications, whereas predictive output entropy exhibits the inverse behavior \cite{hajji2025map}.
\end{enumerate}

Furthermore, a persistent disconnect exists between detection engines and downstream intervention: most systems operate as post-hoc diagnostic monitors rather than active mitigation controllers \cite{dhuliawala2024chain}.

To resolve these limitations, we introduce \textbf{TRUEGUARD}, a lightweight, real-time framework that unifies multi-signal intrinsic uncertainty tracking, learned signal fusion, token-level explainability, and active closed-loop mitigation within a single forward pass. TRUEGUARD concurrently extracts four orthogonal signals directly from the model's internal tensor graph:
\begin{itemize}
    \item \textbf{Internal State Probes (ISP):} Evaluates representation vectors across middle-to-late transformer layers using calibrated linear probes.
    \item \textbf{Attention Anomaly Score (AAS):} Monitors attention weight dispersal over prompt context tokens using normalized Shannon entropy.
    \item \textbf{Entropy-Energy Monitor (EEM):} Computes a weighted combination of output distribution Shannon entropy and R\'enyi-2 confidence.
    \item \textbf{Distribution Shift Tracker (DST):} Measures temporal cosine divergence between successive hidden state vectors during generation.
\end{itemize}

These signals are standardized on-the-fly and integrated via a learned logistic regression fusion head to generate a token-level hallucination risk score $R(t)$. If $R(t)$ exceeds a calibrated safety threshold ($\tau = 0.50$), TRUEGUARD interrupts generation mid-sentence to execute active mitigation via targeted Retrieval-Augmented Grounding (Tier 1) or calibrated abstention (Tier 2).

\subsection{Key Contributions}
The main contributions of this work are summarized as follows:
\begin{enumerate}
    \item \textbf{Unified Multi-Signal Architecture:} We formalize and implement a single-pass extraction pipeline integrating ISP, AAS, EEM, and DST, capturing complementary representation, attention, predictive, and temporal signals.
    \item \textbf{Learned Real-Time Fusion:} We show that a calibrated logistic regression layer operating on online $Z$-score standardized signals achieves superior detection accuracy over individual signal baselines with negligible computational overhead ($<0.1$~ms).
    \item \textbf{Rigorous Benchmark Validation:} Evaluated across five benchmark distributions (TruthfulQA, HaluEval-QA, HaluEval-Summ, HaluEval-Dial, and HELM) using 5-fold cross-validation, TRUEGUARD achieves an average AUROC of $0.8821 \pm 0.0170$ and an ablation-verified detection AUROC of $0.8475 \pm 0.0324$.
    \item \textbf{High Efficiency with 11.5$\times$ Speedup:} On standard inference hardware, TRUEGUARD incurs an overhead of only $12.4 \pm 0.5$~ms/token, delivering an 11.5$\times$ speedup compared to multi-sample Semantic Entropy ($142.3$~ms/token).
    \item \textbf{Closed-Loop Mitigation Verification:} In a double-blind human evaluation study ($n = 300$, Cohen's $\kappa = 0.82$), TRUEGUARD's two-tier mitigation reduces hallucination rates from 31.2\% down to 11.3\%---a 63.8\% relative drop---while preserving 88.0\% answer coverage.
    \item \textbf{Complete Reproducibility:} We provide an open-source repository including local CPU smoke-test harnesses, full GPU reproduction scripts, Google Colab notebooks, and pre-computed benchmark metrics.
\end{enumerate}

\section{Related Work}
\subsection{Hallucination Categorization and Taxonomies}
Hallucinations in natural language generation are systematically classified along two major axes: \textit{intrinsic} and \textit{extrinsic} \cite{ji2023survey, alansari2024survey}. Intrinsic hallucinations directly contradict the provided source context or user prompt (e.g., misstating an entity attribute present in the context). Extrinsic hallucinations introduce unverifiable or fabricated assertions absent from both the context and global world knowledge. Huang et al. \cite{huang2024trustllm} established TrustLLM, demonstrating that hallucination frequency correlates inversely with model parameter scale, yet remains pervasive even in state-of-the-art architectures.

\subsection{White-Box Activation and Uncertainty Probing}
Recent research demonstrates that an LLM's internal activations contain latent representations indicative of truthfulness and confidence \cite{azaria2023internal}. Su et al. \cite{su2024mind} introduced MIND, an unsupervised probing technique operating on internal hidden states. Zhang et al. \cite{zhang2025mhad} presented MHAD, showing that truthfulness awareness is concentrated in specific intermediate transformer layers. Concurrently, Kossen et al. \cite{kossen2024semantic} proposed Semantic Entropy Probes (SEPs), which estimate response semantic variance directly from hidden states in a single forward pass.

In the predictive domain, Ma et al. \cite{ma2025semantic} formulated Semantic Energy, measuring raw logit confidence through a Boltzmann-inspired formulation. Moslonka et al. \cite{moslonka2025learned} utilized token-level entropy production rates for black-box detection. In the structural domain, Dasgupta et al. \cite{dasgupta2025hallushift} demonstrated that semantic drift manifests as a measurable cosine distance shift between consecutive layer-wise representations (HaluShift). In attention analysis, Hajji et al. \cite{hajji2025map} mapped hallucination traces through attention entropy, demonstrating that attention over input tokens is especially sensitive to intrinsic hallucinations.

\subsection{Inference-Time Mitigation}
Addressing hallucinations after full sequence generation introduces substantial latency. Chain-of-Verification \cite{dhuliawala2024chain} and RARR \cite{gao2023rarr} generate auxiliary verification queries, multiplying overall inference time. Retrieval-Augmented Generation (RAG) \cite{lewis2020retrieval} and ExpandR \cite{yao2025expandr} ground generation in external knowledge bases, but universal retrieval on every query introduces redundant processing. In contrast, calibrated abstention \cite{wu2025mitigating} trains models to recognize empirical uncertainty bounds and decline answering when confidence is low. TRUEGUARD synthesizes these insights by deploying targeted RAG and calibrated abstention conditionally based on real-time token risk.

\section{Problem Formulation}
Let an autoregressive language model $\mathcal{M}$ parameterized by $\theta$ generate a sequence of output tokens $Y = (y_1, y_2, \dots, y_T)$ given an input prompt context $C = (c_1, c_2, \dots, c_M)$. At generation step $t \in \{1, \dots, T\}$, the model performs a forward pass:
\begin{equation}
    p(y_t \mid C, y_{<t}) = \text{softmax}(z_t)
\end{equation}
where $z_t \in \mathbb{R}^{|V|}$ denotes the output logit vector over vocabulary $V$. Concurrent with logit computation, the model produces residual hidden representations $\{h_t^L\}_{L=1}^{N_L}$ and multi-head attention matrices $\{A_t^L\}_{L=1}^{N_L}$, where $N_L$ is the total number of transformer layers.

The objective of real-time hallucination detection is to learn an inline mapping:
\begin{equation}
    f: (z_t, \{h_t^L\}, \{A_t^L\}) \mapsto R(t) \in [0, 1]
\end{equation}
such that $R(t)$ accurately reflects the probability that token $y_t$ initiates or constitutes a factual hallucination, evaluated with minimal latency overhead $\Delta t \ll t_{\text{forward}}$.

\begin{figure}[t]
\centering
\includegraphics[width=\columnwidth]{fig1_architecture.png}
\caption{Overall architecture of the TRUEGUARD framework. Four complementary intrinsic signals (ISP, AAS, EEM, DST) are extracted concurrently in a single forward pass, standardized online, and fused via logistic regression to trigger active mitigation.}
\label{fig:architecture}
\end{figure}

\section{The TRUEGUARD Framework}
\subsection{Overall Architecture}
As illustrated in Fig.~\ref{fig:architecture}, TRUEGUARD operates as an inline tensor monitor that intercepts the model's internal activations at each autoregressive step. Rather than invoking secondary forward passes or external verification calls, TRUEGUARD extracts a 4-dimensional feature vector $X_t = [\text{ISP}_t, \text{AAS}_t, \text{EEM}_t, \text{DST}_t]^T$ directly from the ongoing forward execution graph.

\subsection{Signal 1: Internal State Probes (ISP)}
Following representations established by MHAD \cite{zhang2025mhad} and MIND \cite{su2024mind}, truthfulness-related information is concentrated in middle-to-late transformer layers. For a 32-layer model (e.g., LLaMA-2-7B), we select probe layers $L_{\text{isp}} = \{15, 20, 25, 30\}$ (or relative depth percentiles $\{0.45, 0.60, 0.75, 0.90\}$).

For each selected layer $L \in L_{\text{isp}}$, a linear classification probe with weights $w_L \in \mathbb{R}^{d_h}$ and bias $b_L \in \mathbb{R}$ projects the layer hidden state $h_t^L \in \mathbb{R}^{d_h}$ to a scalar probability:
\begin{equation}
    \text{ISP}_t^L = \sigma(w_L^T h_t^L + b_L)
\end{equation}
The aggregated ISP score at step $t$ is computed as the arithmetic mean across probed layers:
\begin{equation}
    \text{ISP}_t = \frac{1}{|L_{\text{isp}}|} \sum_{L \in L_{\text{isp}}} \text{ISP}_t^L
\end{equation}
Each probe is trained using binary cross-entropy loss with $L_2$ weight regularization ($\lambda = 10^{-4}$).

\subsection{Signal 2: Attention Anomaly Score (AAS)}
When generating factually accurate tokens, transformer attention heads attend predominantly to relevant context tokens in $C$. Conversely, during hallucinated generation, attention weights frequently disperse across irrelevant previous tokens \cite{hajji2025map}.

We quantify attention dispersal by computing the normalized Shannon entropy of attention matrices in the final quartile of layers $L_{\text{aas}} = \{25, 26, \dots, 31\}$. For head $h \in \{1, \dots, |H|\}$ at layer $L$, the attention weight distribution from token position $t$ to all prior tokens $i \in \{1, \dots, t\}$ is denoted $A_{t, i}^{h, L}$. The normalized attention entropy is:
\begin{equation}
    H^{h, L}(t) = -\sum_{i=1}^t A_{t, i}^{h, L} \log(A_{t, i}^{h, L} + \epsilon)
\end{equation}
\begin{equation}
    H_{\text{norm}}^{h, L}(t) = \frac{H^{h, L}(t)}{\log(t)}
\end{equation}
where $\epsilon = 10^{-10}$ ensures numerical stability. The composite AAS score averages across all heads and monitored layers:
\begin{equation}
    \text{AAS}_t = \frac{1}{|L_{\text{aas}}| \cdot |H|} \sum_{L \in L_{\text{aas}}} \sum_{h=1}^{|H|} H_{\text{norm}}^{h, L}(t)
\end{equation}
where $\text{AAS}_t \in [0, 1]$. Values approaching 1 denote dispersed attention patterns indicative of ungrounded token generation.

\subsection{Signal 3: Entropy-Energy Monitor (EEM)}
Output token uncertainty is captured by examining the unnormalized logit vector $z_t \in \mathbb{R}^{|V|}$ and the resulting softmax distribution $p_t = \text{softmax}(z_t)$. Standard Shannon entropy measures overall predictive dispersion, but is bounded by softmax temperature scaling. We augment Shannon entropy with R\'enyi-2 entropy, which specifically quantifies probability concentration:
\begin{equation}
    H_{\text{norm}}(t) = \frac{-\sum_{k=1}^{|V|} p_{t, k} \log(p_{t, k} + \epsilon)}{\log(|V|)}
\end{equation}
\begin{equation}
    E_{\text{norm}}(t) = \frac{-\log\left(\sum_{k=1}^{|V|} p_{t, k}^2 + \epsilon\right)}{\log(|V|)}
\end{equation}
The unified EEM score is a convex combination tuned on validation data:
\begin{equation}
    \text{EEM}_t = w_e \cdot H_{\text{norm}}(t) + (1 - w_e) \cdot E_{\text{norm}}(t)
\end{equation}
where $w_e = 0.6$ balances distribution tail uncertainty with mode sharpness.

\subsection{Signal 4: Distribution Shift Tracker (DST)}
Drawing on the cognitive drift hypothesis formalized in HaluShift \cite{dasgupta2025hallushift}, factual hallucination often manifests as a progressive geometric displacement in activation space. Let $\bar{h}_t$ represent the average hidden state across deep layers $L_{\text{dst}} = \{20, 25, 30\}$:
\begin{equation}
    \bar{h}_t = \frac{1}{|L_{\text{dst}}|} \sum_{L \in L_{\text{dst}}} h_t^L
\end{equation}
The instantaneous distribution shift is measured as the normalized cosine distance relative to the state at step $t-1$:
\begin{equation}
    d_t = \frac{1 - \cos(\bar{h}_t, \bar{h}_{t-1})}{2}
\end{equation}
with $d_1 = 0$. The running DST score accumulates this displacement:
\begin{equation}
    \text{DST}_t = \frac{1}{t-1} \sum_{t'=2}^t d_{t'}
\end{equation}
Sudden trajectory deviations elevate $\text{DST}_t$, providing a temporal continuity signal that complements static single-step metrics.

\subsection{Feature Normalization and Learned Fusion}
Because the four signals operate across distinct numerical scales, raw concatenation degrades gradient stability. TRUEGUARD applies an online Welford $Z$-score scaler:
\begin{equation}
    \hat{X}_t = \frac{X_t - \mu_t}{\sigma_t + \epsilon_s}
\end{equation}
where $\mu_t$ and $\sigma_t$ denote running empirical mean and standard deviation vectors updated at each step.

The standardized vector $\hat{X}_t \in \mathbb{R}^4$ is mapped to a unified risk score $R(t)$ via learned logistic regression:
\begin{equation}
    R(t) = \sigma(W^T \hat{X}_t + b)
\end{equation}
where $W = [w_{\text{isp}}, w_{\text{aas}}, w_{\text{eem}}, w_{\text{dst}}]^T$ and $b \in \mathbb{R}$ are optimized using binary cross-entropy on validation benchmarks.

To evaluate full response sequences, the sequence risk score incorporates a frequency penalty for persistent high-risk tokens:
\begin{equation}
    R_{\text{seq}} = \min\left(1.0, \, \max_t R(t) \cdot \left[1 + \gamma \cdot \frac{\sum_{t=1}^T \mathbb{I}(R(t) > \tau)}{T}\right]\right)
\label{eq:rseq}
\end{equation}
where $\tau = 0.50$ is the decision threshold and $\gamma = 0.20$ penalizes multi-token hallucination cascades.

\subsection{Active Closed-Loop Mitigation}
When $R(t) > \tau$, TRUEGUARD immediately halts standard generation mid-sentence and initiates a two-tier mitigation response:
\begin{itemize}
    \item \textbf{Tier 1 (Targeted RAG):} Key entities from the prompt context are extracted to query an external factual knowledge store (e.g., Wikipedia API). Retrieved passages are injected into an augmented prompt, and generation is restarted with factual grounding.
    \item \textbf{Tier 2 (Calibrated Abstention):} If retrieval fails to locate high-confidence evidence or if the regenerated sequence risk remains above $\tau$, TRUEGUARD executes calibrated abstention, emitting an epistemically honest disclaimer: \textit{``I do not have sufficient verified evidence to answer this query reliably without risk of factual inaccuracy.''}
\end{itemize}

\section{Experimental Methodology}
\subsection{Models and Hardware Environment}
We evaluate TRUEGUARD across three primary causal language model families:
\begin{itemize}
    \item \textbf{Primary Evaluation Models:} LLaMA-2-7B-Chat \cite{touvron2023llama} and Mistral-7B-Instruct-v0.2 \cite{jiang2023mistral} (32 transformer layers, 4096 hidden dimension, 32 attention heads).
    \item \textbf{Lightweight Validation Model:} TinyLlama-1.1B-Chat \cite{zhang2024tinyllama} (22 layers, 2048 hidden dimension, 32 heads).
    \item \textbf{CPU Smoke-Test Baseline:} GPT-2 (12 layers, 768 hidden dimension, 12 heads).
\end{itemize}

Inference experiments were conducted in both low-cost cloud environments (NVIDIA T4 GPU with 16~GB VRAM) and local standard CPU workstations to verify cross-platform portability.

\subsection{Evaluation Benchmarks}
Our evaluation encompasses five standardized hallucination detection benchmarks:
\begin{enumerate}
    \item \textbf{TruthfulQA} \cite{lin2022truthfulqa}: 817 questions spanning 38 categories designed to induce common misconceptions and false beliefs.
    \item \textbf{HaluEval-QA} \cite{li2023halueval}: 10,000 question-answer pairs with paired factual and hallucinated responses.
    \item \textbf{HaluEval-Summ}: Knowledge-grounded text summarization with synthetic entity fabrications.
    \item \textbf{HaluEval-Dial}: Dialogue completion evaluating conversational consistency.
    \item \textbf{HELM (Holistic Evaluation of Language Models)} \cite{liang2022holistic}: Open-domain knowledge retrieval and multi-step reasoning queries.
\end{enumerate}

\subsection{Baselines}
We benchmark TRUEGUARD against leading representative methods across black-box, single-signal white-box, and multi-signal categories:
\begin{itemize}
    \item \textbf{Token Entropy (Black-Box):} Baseline predictive entropy of the output distribution.
    \item \textbf{Semantic Entropy \cite{kossen2024semantic}:} Multi-sample clustering baseline sampling 5 to 10 stochastic completions per prompt.
    \item \textbf{MIND \cite{su2024mind}:} Unsupervised probing of internal layer hidden states.
    \item \textbf{LLM-Check \cite{sriramanan2024llmcheck}:} Single-pass hidden state classifier.
    \item \textbf{HaluShift \cite{dasgupta2025hallushift}:} Single-signal distribution shift tracker monitoring cosine state distance.
\end{itemize}

\subsection{Evaluation Metrics}
Detection efficacy is evaluated using Area Under the Receiver Operating Characteristic Curve (AUROC), Area Under the Precision-Recall Curve (AUPRC), F1 Score, Precision, Recall, and False Positive Rate (FPR). Every AUROC value is reported via 5-fold cross-validation, including sample standard deviation and 95\% confidence intervals:
\begin{equation}
    \text{CI}_{95\%} = \bar{x} \pm 1.96 \cdot \frac{s}{\sqrt{k}}
\end{equation}
where $k = 5$ denotes cross-validation folds.

Mitigation efficacy is quantified by relative hallucination reduction:
\begin{equation}
    \text{Reduction (\%)} = \frac{\text{Rate}_{\text{before}} - \text{Rate}_{\text{after}}}{\text{Rate}_{\text{before}}} \times 100\%
\end{equation}
and answer coverage percentage (fraction of queries answered without abstaining).

\section{Results and Empirical Analysis}

\subsection{Overall Benchmark Performance}
Table~\ref{tab:benchmark_auroc} presents 5-fold cross-validation performance across all five benchmarks. TRUEGUARD achieves consistent detection accuracy, obtaining an overall average AUROC of $0.8821 \pm 0.0170$. On TruthfulQA, TRUEGUARD achieves an AUROC of $0.8533 \pm 0.0533$ (95\% CI: $[0.8066, 0.9001]$), significantly outperforming Token Entropy ($0.712$) and Semantic Entropy ($0.776$).

\begin{table}[t]
\caption{Cross-Validation Performance Across Benchmarks (5-Fold CV)}
\label{tab:benchmark_auroc}
\centering
\begin{tabular}{lcccc}
\toprule
\textbf{Benchmark} & \textbf{AUROC} & \textbf{95\% CI} & \textbf{AUPRC} & \textbf{F1 Score} \\
\midrule
TruthfulQA       & 0.8533 $\pm$ 0.0533 & [0.8066, 0.9001] & 0.8140 & 0.7186 \\
HaluEval-QA      & 0.8840 $\pm$ 0.0185 & [0.8678, 0.9003] & 0.8739 & 0.7874 \\
HaluEval-Summ    & 0.8797 $\pm$ 0.0440 & [0.8411, 0.9184] & 0.8459 & 0.7252 \\
HaluEval-Dial    & 0.8872 $\pm$ 0.0387 & [0.8533, 0.9212] & 0.8484 & 0.7373 \\
HELM             & 0.9063 $\pm$ 0.0454 & [0.8665, 0.9461] & 0.8552 & 0.7545 \\
\midrule
\textbf{Overall Mean} & \textbf{0.8821 $\pm$ 0.0170} & --- & \textbf{0.8475} & \textbf{0.7446} \\
\bottomrule
\end{tabular}
\end{table}

Fig.~\ref{fig:auroc_comparison} compares TRUEGUARD against baseline methods. TRUEGUARD achieves the highest AUROC across all task categories, demonstrating that multi-signal fusion overcomes the individual blind spots of single-signal detectors.

\begin{figure}[t]
\centering
\includegraphics[width=\columnwidth]{fig2_auroc_benchmark_comparison.png}
\caption{Hallucination detection AUROC comparison across benchmarks. TRUEGUARD consistently outperforms black-box and single-signal white-box methods.}
\label{fig:auroc_comparison}
\end{figure}

\subsection{Ablation Study: Contribution of Individual Signals}
To verify the hypothesis that all four intrinsic signals provide complementary diagnostic value, we conducted a systematic ablation experiment on the held-out validation suite. Table~\ref{tab:ablation} details the detection AUROC and performance degradation ($\Delta$) when removing each signal.

\begin{table}[t]
\caption{Ablation Study: Signal Contribution to AUROC}
\label{tab:ablation}
\centering
\begin{tabular}{lccc}
\toprule
\textbf{Configuration} & \textbf{AUROC} & \textbf{Std} & $\mathbf{\Delta}$ \textbf{from Full} \\
\midrule
\textbf{Full TRUEGUARD (4 Signals)} & \textbf{0.8475} & $\pm$ 0.0324 & \textbf{Reference} \\
$- \text{EEM}$ (No Entropy/Energy)  & 0.7957 & $\pm$ 0.0316 & $-0.0518$ \\
$- \text{ISP}$ (No Hidden Probes)   & 0.8093 & $\pm$ 0.0465 & $-0.0382$ \\
$- \text{DST}$ (No Dist. Shift)     & 0.8320 & $\pm$ 0.0237 & $-0.0155$ \\
$- \text{AAS}$ (No Attention)       & 0.8370 & $\pm$ 0.0395 & $-0.0105$ \\
\midrule
EEM Only                            & 0.7516 & $\pm$ 0.0395 & $-0.0959$ \\
ISP Only                            & 0.7445 & $\pm$ 0.0537 & $-0.1030$ \\
DST Only                            & 0.6812 & $\pm$ 0.0626 & $-0.1663$ \\
AAS Only                            & 0.6317 & $\pm$ 0.0345 & $-0.2158$ \\
\bottomrule
\end{tabular}
\end{table}

\begin{figure}[t]
\centering
\includegraphics[width=\columnwidth]{fig3_ablation_study.png}
\caption{Ablation study illustrating performance degradation when omitting individual signals. Removing EEM causes the sharpest drop ($-0.0518$), followed by ISP ($-0.0382$).}
\label{fig:ablation}
\end{figure}

As shown in Fig.~\ref{fig:ablation}, removing the Entropy-Energy Monitor produces the largest single degradation ($\Delta = -0.0518$), indicating that predictive probability distributions provide the primary signal for factual uncertainty. Removing Internal State Probes degrades AUROC by $-0.0382$, confirming that intermediate hidden representations encode semantic truthfulness cues that logits alone cannot capture. Importantly, every single-signal model operates substantially below the full 4-signal fusion (e.g., EEM only achieves $0.7516$ vs.\ $0.8475$ for Full TRUEGUARD), proving mutual complementarity.

\subsection{Latency and Computational Efficiency}
Table~\ref{tab:latency} details the empirical execution time breakdown measured per token on an NVIDIA T4 GPU and standard CPU testbeds.

\begin{table}[t]
\caption{Inference Latency and Computational Overhead}
\label{tab:latency}
\centering
\begin{tabular}{lccc}
\toprule
\textbf{Method} & \textbf{Passes/Tok} & \textbf{Overhead (ms)} & \textbf{Speedup vs SE} \\
\midrule
Semantic Entropy \cite{kossen2024semantic} & 5--10 & 142.3 $\pm$ 3.8 & 1.0$\times$ (Baseline) \\
LLM-Check \cite{sriramanan2024llmcheck}    & 1     & 31.5 $\pm$ 1.2  & 4.5$\times$ \\
HaluShift \cite{dasgupta2025hallushift}    & 1     & 18.2 $\pm$ 0.6  & 7.8$\times$ \\
\textbf{TRUEGUARD (Ours)}                  & \textbf{1} & \textbf{12.4 $\pm$ 0.5} & \textbf{11.5$\times$} \\
\bottomrule
\end{tabular}
\end{table}

\begin{figure}[t]
\centering
\includegraphics[width=\columnwidth]{fig4_latency_efficiency_tradeoff.png}
\caption{Detection AUROC vs. Latency overhead (log scale). TRUEGUARD resides in the optimal low-latency, high-accuracy quadrant.}
\label{fig:latency}
\end{figure}

While Semantic Entropy demands $142.3 \pm 3.8$~ms/token due to multi-sample generation loops, TRUEGUARD executes in $12.4 \pm 0.5$~ms/token---achieving an \textbf{11.5$\times$ speedup}. Component profiling confirms that logistic regression fusion requires $<0.1$~ms, with signal extraction requiring $0.49$~ms (ISP), $1.26$~ms (AAS), $0.47$~ms (EEM), and $0.26$~ms (DST). Fig.~\ref{fig:latency} visualizes this efficiency frontier, demonstrating TRUEGUARD's position in the optimal quadrant.

\subsection{Mitigation Effectiveness}
We evaluated the two-tier mitigation pipeline using double-blind human annotation on $n = 300$ model responses sampled evenly across TruthfulQA, HaluEval, and HELM. Two domain experts evaluated factual faithfulness independently, achieving high inter-rater agreement (Cohen's $\kappa = 0.82$). Discrepancies were resolved via consensus adjudication.

\begin{table}[t]
\caption{Mitigation Pipeline Impact on Hallucination Rate ($n=300$)}
\label{tab:mitigation}
\centering
\begin{tabular}{lccc}
\toprule
\textbf{Mitigation Stage} & \textbf{Halluc. Rate} & \textbf{Accuracy} & \textbf{Coverage} \\
\midrule
Base LLM (No Guard)        & 31.2\% & 68.8\% & 100.0\% \\
Detection Flag Only         & 31.2\% & 68.8\% & 100.0\% \\
+ Targeted RAG (Tier 1)     & 18.4\% & 79.3\% & 100.0\% \\
+ Calibrated Abstention (Tier 2) & 12.7\% & 84.1\% & 95.0\% \\
\textbf{Full TRUEGUARD Closed-Loop} & \textbf{11.3\%} & \textbf{85.6\%} & \textbf{88.0\%} \\
\bottomrule
\end{tabular}
\end{table}

\begin{figure}[t]
\centering
\includegraphics[width=\columnwidth]{fig5_mitigation_waterfall.png}
\caption{Reduction in hallucination rate across mitigation stages ($n=300, \kappa=0.82$). TRUEGUARD achieves a 63.8\% relative reduction.}
\label{fig:mitigation}
\end{figure}

As presented in Table~\ref{tab:mitigation} and Fig.~\ref{fig:mitigation}, the baseline hallucination rate of 31.2\% is reduced to 18.4\% through Tier-1 Targeted RAG. Adding Tier-2 Calibrated Abstention drives the final hallucination rate down to 11.3\%, yielding a \textbf{63.8\% relative reduction}:
\begin{equation}
    \text{Relative Reduction} = \frac{31.2\% - 11.3\%}{31.2\%} = \frac{19.9\%}{31.2\%} = 63.78\% \approx 63.8\%
\end{equation}
Importantly, TRUEGUARD achieves this reduction while maintaining an 88.0\% query answer coverage, abstaining on only 12.0\% of queries where verified facts could not be retrieved.

\subsection{Token-Level Risk Visualization and Error Analysis}
Fig.~\ref{fig:token_trace} illustrates a real-time risk trace during generation of an induced hallucination prompt (\textit{``In 1965, the first person to walk on Mars was Soviet cosmonaut\dots''}). 

\begin{figure}[t]
\centering
\includegraphics[width=\columnwidth]{fig6_token_risk_heatmap.png}
\caption{Real-time token-level risk trace during hallucinated sequence generation. $R(t)$ surges above safety threshold $\tau=0.50$ at the fabricated entity, triggering mid-sentence interception.}
\label{fig:token_trace}
\end{figure}

During initial context processing, all four signals remain low ($R(t) \approx 0.11$--$0.21$). Upon generating the fictitious entity \textit{``Alexei Leonov''}, predictive uncertainty (EEM) spikes to $0.88$, attention entropy (AAS) increases to $0.74$, and the fused risk score reaches $0.83 > \tau$. TRUEGUARD halts generation at token step 6, intercepting the hallucinated cascade before text emission completes.

\section{Discussion}
\subsection{Complementarity of Detection Signals}
Our experimental findings confirm the theoretical hypothesis that hallucination in autoregressive language models is not a monolithic phenomenon. Instead:
\begin{itemize}
    \item \textbf{EEM} serves as the primary detector for output vocabulary uncertainty;
    \item \textbf{ISP} identifies latent truthfulness discrepancies in intermediate representations;
    \item \textbf{AAS} isolates instances where the model loses focus on user-supplied constraints;
    \item \textbf{DST} tracks progressive temporal deviation from the factual trajectory.
\end{itemize}
Because each signal targets an orthogonal mathematical subspace of the transformer forward pass, linear fusion provides robust cross-domain generalization.

\subsection{Ethical Considerations}
Deploying active mitigation mechanisms introduces epistemic trade-offs. While calibrated abstention reduces false assertions, overly conservative thresholds risk reducing system utility. Transparency in flagging whether an output was grounded via Tier-1 RAG or altered via Tier-2 Abstention is essential to prevent user deception and preserve appropriate trust calibration.

\section{Limitations}
Several limitations should be noted:
\begin{enumerate}
    \item \textbf{White-Box Model Access:} TRUEGUARD relies on access to hidden states, attention matrices, and logits. Consequently, it cannot be directly applied to proprietary black-box APIs (e.g., GPT-4) without access to provider log-probabilities.
    \item \textbf{Knowledge Base Quality:} The effectiveness of Tier-1 Targeted RAG depends directly on external database completeness. Ambiguous or missing entries trigger Tier-2 abstention.
    \item \textbf{Language Scope:} Primary validation was performed on English benchmarks. Extending probe calibration to multilingual corpora remains an area for future work.
\end{enumerate}

\section{Reproducibility Statement}
To guarantee complete scientific reproducibility, all project artifacts are open-sourced:
\begin{itemize}
    \item Complete Python package under \texttt{TRUEGUARD\_FINAL/src/};
    \item Automated 60-second CPU smoke-test script (\texttt{experiments/cpu\_smoke\_test.py});
    \item 5-fold cross-validation runner (\texttt{experiments/run\_benchmark.py});
    \item Standalone ablation suite (\texttt{experiments/ablation.py});
    \item Self-contained Google Colab notebook (\texttt{notebooks/TRUEGUARD\_Colab.ipynb});
    \item All raw metrics and generated 300 DPI figures (\texttt{results/}, \texttt{figures/}).
\end{itemize}

\section{Conclusion}
In this paper, we presented TRUEGUARD, a unified, lightweight framework for real-time hallucination detection and mitigation in LLMs. By extracting four complementary intrinsic signals (ISP, AAS, EEM, DST) within a single forward pass and fusing them via calibrated logistic regression, TRUEGUARD achieves an average AUROC of $0.8821 \pm 0.0170$ across standard benchmarks with only $12.4 \pm 0.5$~ms/token overhead---an 11.5$\times$ speedup over multi-sample Semantic Entropy. Furthermore, TRUEGUARD's active closed-loop mitigation achieves a 63.8\% relative reduction in hallucination occurrences while preserving 88.0\% query answer coverage. TRUEGUARD demonstrates that robust, explainable, and active hallucination guardrails can be deployed in production LLM environments without sacrificing real-time generation speed.

\section*{Acknowledgment}
The authors thank the Department of Computer Science and Engineering at Vellore Institute of Technology, Chennai, for providing academic support and computational resources.

\bibliographystyle{IEEEtran}
\begin{thebibliography}{10}

\bibitem{huang2024trustllm}
Y.~Huang, L.~Sun, H.~Wang, S.~Wu, Q.~Zhang \emph{et~al.}, ``TrustLLM: Trustworthiness in large language models,'' \emph{arXiv preprint arXiv:2401.05561}, 2024.

\bibitem{zhang2025siren}
Y.~Zhang, Y.~Li, L.~Cui, D.~Cai, L.~Liu, T.~Fu \emph{et~al.}, ``Siren's song in the AI ocean: A survey on hallucination in large language models,'' \emph{Computational Linguistics}, vol.~51, no.~4, pp. 1373--1418, 2025.

\bibitem{ji2023survey}
Z.~Ji, N.~Lee, R.~Frieske, T.~Yu, D.~Su, Y.~Xu \emph{et~al.}, ``Survey of hallucination in natural language generation,'' \emph{ACM Computing Surveys}, vol.~55, no.~12, pp. 1--38, 2023.

\bibitem{min2023factscore}
S.~Min, K.~Krishna, X.~Lyu, M.~Lewis, W.-t. Yih, P.~W. Koh \emph{et~al.}, ``FActScore: Fine-grained atomic evaluation of factual precision in long form text generation,'' in \emph{Proc. Conf. Empirical Methods Natural Lang. Process. (EMNLP)}, 2023, pp. 12076--12100.

\bibitem{manakul2023selfcheckgpt}
P.~Manakul, A.~Liusie, and M.~Gales, ``SelfCheckGPT: Zero-resource black-box hallucination detection for generative large language models,'' in \emph{Proc. Conf. Empirical Methods Natural Lang. Process. (EMNLP)}, 2023, pp. 9004--9017.

\bibitem{kossen2024semantic}
J.~Kossen, J.~Han, M.~Razzak, L.~Schut, S.~Malik, and Y.~Gal, ``Semantic entropy probes: Robust and cheap hallucination detection in LLMs,'' in \emph{ICML Workshop on Foundation Models in the Wild}, 2024.

\bibitem{su2024mind}
W.~Su, C.~Wang, Q.~Ai, Y.~Hu, Z.~Wu, Y.~Zhou, and Y.~Liu, ``Unsupervised real-time hallucination detection based on the internal states of large language models,'' in \emph{Findings Assoc. Comput. Linguistics: ACL}, 2024, pp. 14379--14391.

\bibitem{zhang2025mhad}
L.~Zhang, D.~Song, Z.~Wu, Y.~Tian, C.~Zhou, J.~Xu \emph{et~al.}, ``Detecting hallucination in large language models through deep internal representation analysis,'' in \emph{Proc. Int. Joint Conf. Artif. Intell. (IJCAI)}, 2025, pp. 8357--8365.

\bibitem{hajji2025map}
E.~Hajji, A.~Bouguerra, and F.~Arnez, ``The map of misbelief: Tracing intrinsic and extrinsic hallucinations through attention patterns,'' in \emph{Proc. AAAI Symp. Ser.}, vol.~7, no.~1, 2025, pp. 177--184.

\bibitem{ma2025semantic}
H.~Ma, J.~Pan, J.~Liu, Y.~Chen, J.~T. Zhou, G.~Wang \emph{et~al.}, ``Semantic energy: Detecting LLM hallucination beyond entropy,'' \emph{arXiv preprint arXiv:2508.14496}, 2025.

\bibitem{dasgupta2025hallushift}
S.~Dasgupta, S.~Nath, A.~Basu, P.~Shamsolmoali, and S.~Das, ``HaluShift: Measuring distribution shifts towards hallucination detection in LLMs,'' in \emph{Proc. IEEE Int. Joint Conf. Neural Netw. (IJCNN)}, 2025, pp. 1--8.

\bibitem{dhuliawala2024chain}
S.~Dhuliawala, M.~Komeili, J.~Xu, R.~Raileanu, X.~Li, A.~Celikyilmaz, and J.~Weston, ``Chain-of-verification reduces hallucination in large language models,'' in \emph{Findings Assoc. Comput. Linguistics: ACL}, 2024, pp. 3563--3578.

\bibitem{alansari2024survey}
A.~Alansari and H.~Luqman, ``Large language models hallucination: A comprehensive survey,'' \emph{arXiv preprint arXiv:2404.00806}, 2024.

\bibitem{azaria2023internal}
A.~Azaria and T.~Mitchell, ``The internal state of an LLM knows when it's lying,'' in \emph{Findings Assoc. Comput. Linguistics: EMNLP}, 2023, pp. 967--976.

\bibitem{moslonka2025learned}
C.~Moslonka, H.~Randrianarivo, A.~Garnier, and E.~Malherbe, ``Learned hallucination detection in black-box LLMs using token-level entropy production rate,'' \emph{arXiv preprint arXiv:2509.04492}, 2025.

\bibitem{gao2023rarr}
L.~Gao, Z.~Dai, P.~Pasupat, A.~Chen, A.~T. Chaganty, Y.~Fan \emph{et~al.}, ``RARR: Researching and revising what language models say, using language models,'' in \emph{Proc. 61st Annu. Meeting Assoc. Comput. Linguistics (ACL)}, 2023, pp. 16477--16508.

\bibitem{lewis2020retrieval}
P.~Lewis, E.~Perez, A.~Piktus, F.~Petroni, V.~Karpukhin \emph{et~al.}, ``Retrieval-augmented generation for knowledge-intensive NLP tasks,'' in \emph{Advances in Neural Inf. Process. Syst. (NeurIPS)}, vol.~33, 2020, pp. 9459--9474.

\bibitem{yao2025expandr}
S.~Yao, P.~Huang, Z.~Liu, Y.~Gu, Y.~Yan, S.~Yu, and G.~Yu, ``ExpandR: Teaching dense retrievers beyond queries with LLM guidance,'' in \emph{Proc. Conf. Empirical Methods Natural Lang. Process. (EMNLP)}, 2025, pp. 19036--19054.

\bibitem{wu2025mitigating}
J.~Wu, J.~Liu, Z.~Zeng, T.~Zhan, T.~Cai, and W.~Huang, ``Mitigating LLM hallucination via behaviorally calibrated reinforcement learning,'' \emph{arXiv preprint arXiv:2512.19920}, 2025.

\bibitem{touvron2023llama}
H.~Touvron, L.~Martin, K.~Stone, P.~Albert, A.~Almahairi \emph{et~al.}, ``Llama 2: Open foundation and fine-tuned chat models,'' \emph{arXiv preprint arXiv:2307.09288}, 2023.

\bibitem{jiang2023mistral}
A.~Q. Jiang, A.~Sablayrolles, A.~Mensch, C.~Bamford, D.~S. Chaplot \emph{et~al.}, ``Mistral 7B,'' \emph{arXiv preprint arXiv:2310.06825}, 2023.

\bibitem{zhang2024tinyllama}
P.~Zhang, G.~Zeng, T.~Wang, and W.~Lu, ``TinyLlama: An open-source small language model,'' \emph{arXiv preprint arXiv:2401.02385}, 2024.

\bibitem{lin2022truthfulqa}
S.~Lin, J.~Hilton, and O.~Evans, ``TruthfulQA: Measuring how models mimic human falsehoods,'' in \emph{Proc. 60th Annu. Meeting Assoc. Comput. Linguistics (ACL)}, 2022, pp. 3214--3252.

\bibitem{li2023halueval}
J.~Li, X.~Cheng, W.~X. Zhao, J.-Y. Nie, and J.-R. Wen, ``HaluEval: A large-scale hallucination evaluation benchmark for large language models,'' in \emph{Proc. Conf. Empirical Methods Natural Lang. Process. (EMNLP)}, 2023, pp. 6449--6464.

\bibitem{liang2022holistic}
P.~Liang, R.~Bommasani, T.~Lee, D.~Tsipras, D.~Soylu \emph{et~al.}, ``Holistic evaluation of language models,'' \emph{arXiv preprint arXiv:2211.09110}, 2022.

\bibitem{sriramanan2024llmcheck}
G.~Sriramanan, S.~Bharti, V.~S. Sadasivan, S.~Saha, P.~Kattakinda, and S.~Feizi, ``LLM-Check: Investigating detection of hallucinations in large language models,'' in \emph{Advances in Neural Inf. Process. Syst. (NeurIPS)}, vol.~37, 2024, pp. 34188--34216.

\end{thebibliography}

\end{document}
